% payoff.tex — Fee concentration index, null hypothesis, deviation, price mapping
% Kontrol: prove_cfmm_feeconc_indexBoundedness, prove_cfmm_feeconc_nullConsistency

\section{Fee Concentration Index}\label{sec:fee-concentration-index}

\subsection{Primary Index: $\AT$}

The fee concentration index $\AT$ is a path-dependent, HHI-weighted measure of adverse competition
in liquidity provision. It is the \emph{primary} computed and tracked quantity:
\begin{equation}\label{eq:concentration-index}
  \AT = \Bigl(\sum_{k} \thetapos_k \cdot (x_k)^2 \Bigr)^{1/2}
\end{equation}
where the sum runs over all positions removed up to observation time $T$, and:
\begin{itemize}
  \item $x_k = \dfrac{\ell_k}{L(p^*)}$
        is the fee share captured by position $k$, where $\ell_k$ is the
        position's liquidity and $L(p^*)$ is the total liquidity in the
        active tick range;
  \item $\thetapos_k = \dfrac{1}{\Delta B_k}$
        is the sophistication weight, where $\Delta B_k = \text{block}(\text{removal}) - \text{block}(\text{creation})$
        is the position lifetime in blocks;
  \item $\thetapos_k = 1$ for JIT positions ($\Delta B_k = 1$ block: same-block add and remove).
\end{itemize}

The index satisfies $\AT \in [0,1]$.

\subsection{Co-Primary State Variables}

The hook tracks three state variables at each liquidity event:

\begin{definition}[Co-Primary State]\label{def:co-primary}
\begin{align}
  \AT &= \Bigl(\sum_{k} \thetapos_k \cdot x_k^2 \Bigr)^{1/2}
    && \text{fee concentration index} \label{eq:at-state}\\
  \ThetaSum &= \sum_{k} \thetapos_k
    && \text{theta sum (aggregate turnover rate)} \label{eq:theta-state}\\
  \PosCount &= |\{k : \text{position}_k \text{ active}\}|
    && \text{active position count} \label{eq:n-state}
\end{align}
All three are stored on-chain. The derived quantities $\ATnull$ and $\DeltaPlus$
are computed on the fly from these three values.
\end{definition}

The Ma--Crapis \cite{macrapis2024permissionless} symmetric equilibrium assumes all positions capture equal fee shares
$x_k = 1/\PosCount$. Under this null hypothesis:
\begin{equation}\label{eq:null-at}
  \ATnull = \sqrt{\frac{\ThetaSum}{\PosCount^2}}
\end{equation}

\begin{proposition}[Null as Lower Bound]\label{prop:null-lower-bound}
For any fee share distribution $\{x_k\}$ with $\sum_k x_k = 1$ and $x_k \geq 0$:
\begin{equation}
  \AT \geq \ATnull
\end{equation}
with equality if and only if $x_k = 1/\PosCount$ for all $k$.
\end{proposition}

\begin{proof}

 Under equal shares $x_k = 1/\PosCount$:
\[
  (\ATnull)^2 = \sum_k \thetapos_k \cdot \frac{1}{\PosCount^2}
             = \frac{\ThetaSum}{\PosCount^2}
\]

For arbitrary shares, by Jensen's inequality on the convex function $x \mapsto x^2$:
\[
  \sum_k \thetapos_k x_k^2 \geq \frac{\ThetaSum}{\PosCount^2}
  \quad \text{when } \sum_k x_k = 1
\]
since $\sum_k x_k^2 \geq 1/\PosCount$ by the constraint $\sum x_k = 1$
(this is the HHI lower bound). Multiplying each term by $\thetapos_k/\ThetaSum$
and using convexity preserves the inequality. Taking square roots gives $\AT \geq \ATnull$. \qedhere
\end{proof}

\subsection{Concentration Deviation: $\DeltaPlus$}\label{subsec:deviation}

The economically meaningful treatment variable is the excess concentration above the
competitive null:

\begin{definition}[Concentration Deviation]\label{def:delta-plus}
\begin{equation}\label{eq:delta-plus}
  \DeltaPlus = \max\bigl(0,\; \AT - \ATnull\bigr)
\end{equation}
\end{definition}

By \cref{prop:null-lower-bound}, $\DeltaPlus \geq 0$ always.
$\DeltaPlus = 0$ corresponds to the competitive equilibrium (all LPs earn equal fee shares).
$\DeltaPlus > 0$ indicates excess concentration: a few LPs capture disproportionate fee revenue.


\subsection{Price Mapping}\label{subsec:price-mapping}

\begin{definition}[Concentration Price]\label{def:price-mapping}
The price of fee concentration in numeraire (USDC) is the odds ratio:
\begin{equation}\label{eq:price-mapping}
  p = \frac{\DeltaPlus}{1 - \DeltaPlus}
\end{equation}
with $p \in [0, \infty)$ for $\DeltaPlus \in [0, 1)$.
The inverse mapping is $\DeltaPlus = p/(1+p)$.
\end{definition}

\begin{proposition}[Monotonicity]\label{prop:price-monotone}
$p(\DeltaPlus)$ is strictly increasing on $[0, 1)$.
\end{proposition}

\begin{proof}
$\dfrac{dp}{d\DeltaPlus} = \dfrac{1}{(1-\DeltaPlus)^2} > 0$ for all $\DeltaPlus < 1$. \qedhere
\end{proof}

\section{Econometric Calibration}\label{sec:econometric}

The econometric identification (see companion document \texttt{lp-insurance-demand.tex})
establishes the demand structure for fee concentration insurance.

\subsection{Quadratic Deviation Model}

The exit hazard model with quadratic treatment:
\begin{equation}\label{eq:quadratic-hazard}
  P(\text{exit}_{i,t} = 1) = \sigma\bigl(
    \beta_0 + \beta_1 \cdot \text{dev}_t + \beta_2 \cdot \text{dev}_t^2
    + \beta_3 \cdot \text{IL}_t + \beta_4 \cdot \ln(\text{age}_{i,t})
  \bigr)
\end{equation}
where $\text{dev}_t = \AT - \ATnull$ is the deviation from the competitive null.

\subsection{Key Results}

\begin{center}
\begin{tabular}{@{}lcccc@{}}
\toprule
Lag & $\beta_1$ (linear) & $p$ & $\beta_2$ (quadratic) & $p$ \\
\midrule
1 & $-23.18$ & $0.012^{**}$ & $+129.20$ & $0.030^{**}$ \\
2 & $-43.42$ & $0.016^{**}$ & $+226.92$ & $0.065^{*}$ \\
3 & $-32.44$ & $0.001^{***}$ & $+205.34$ & $0.004^{***}$ \\
\bottomrule
\end{tabular}
\end{center}

\subsection{Inverted-U and the Insurance Trigger}

The model produces an inverted-U with turning point:
\begin{equation}\label{eq:turning-point}
  \DeltaStar = -\frac{\beta_1}{2\beta_2} \approx 0.09
\end{equation}

\begin{itemize}
  \item \textbf{Below $\DeltaStar$ (shelter regime):} Mild fee concentration correlates with
    high volume. All LPs benefit from elevated fee revenue. $\partial P(\text{exit})/\partial \text{dev} < 0$.
  \item \textbf{Above $\DeltaStar$ (Capponi regime):} Extreme concentration means passive LPs'
    effective fee rate drops below their exit threshold.
    $\partial P(\text{exit})/\partial \text{dev} > 0$.
    Capponi \& Zhu (2024): $\partial p^*_U/\partial \varphi > 0$, so lower $\varphi \to$ exit.
\end{itemize}

The turning point $\DeltaStar \approx 0.09$ identifies when insurance demand activates.
This calibrates the protection mechanism regardless of architecture choice.

\section{Settled Design Properties}\label{sec:design-properties}

The following properties are confirmed by the ETH/USDC 30bps historical backtest
(41 days, 600 positions, 2025-12-05 to 2026-01-14) and hold independent of
implementation details.

\subsection{Backtest-Confirmed Properties}

\begin{enumerate}
  \item \textbf{Perpetual.} No fixed expiry. The instrument settles via a funding rate mechanism
    at liquidity events, following the framework of Angeris et al.\ \cite{angeris2022perpetuals}.

  \item \textbf{Jump-adjusted funding.} $\DeltaPlus$ can jump discontinuously when large positions
    mint or burn, changing $\PosCount$ and $\ThetaSum$ abruptly. The backtest confirms
    concentration jumps are rare but large (1 trigger day in 41 days, $\DeltaPlus$: $0 \to 0.158$).
    The funding rate accounts for this jump risk per the Primer on Perpetuals.

  \item \textbf{Discrete funding at liquidity events.} Funding is exchanged when
    \texttt{afterSwap}, \texttt{afterAddLiquidity}, or \texttt{afterRemoveLiquidity}
    fires in the underlying pool---matching the FCI update frequency.
    No keeper or per-block updates required.

  \item \textbf{USDC numeraire.} Insurance is denominated in USDC (the pool's fee token).

  \item \textbf{Oracle-assisted framework.} The FCI hook provides $\pindex = \DeltaPlus/(1-\DeltaPlus)$
    as the index price. Following Angeris et al.\ \cite{angeris2021replicating}, the oracle-assisted
    replication framework applies because the hook's own state
    ($\AT$, $\ThetaSum$, $\PosCount$) provides the index endogenously.

  \item \textbf{Monotone payoff.} The protection value is monotone non-decreasing in
    $p = \DeltaPlus/(1-\DeltaPlus)$. The inverted-U from econometrics calibrates
    \emph{demand} (where PLPs buy protection), not payoff shape.
\end{enumerate}

\subsection{Optimal Payoff Formula: Power-Squared Exit Payoff}\label{subsec:optimal-payoff}

The backtest compared seven candidate payoff formulas across the ETH/USDC 30bps
position dataset:

\begin{center}
\begin{tabular}{@{}llcc@{}}
\toprule
Mode & Formula & \% Better Off & Mean HV (\$) \\
\midrule
Trigger 1/N (INS-05) & $\text{payout} = \Res / N_{\text{insured}}$ & 0.0\% & $-107.67$ \\
Trigger premium-prop & $\text{payout} \propto \text{premium}_i$ & 2.7\% & $-95.16$ \\
Exit $\times\; \DeltaPlus$ & $\text{payout} = \premfactor \cdot f \cdot \DeltaPlus_{\max}$ & 4.2\% & $-88.03$ \\
Exit $\times\; p$ & $\text{payout} = \premfactor \cdot f \cdot p_{\max}$ & 8.1\% & $-71.22$ \\
\textbf{Exit $\times\; p^2$ ($\alpha=2$)} & $\textbf{payout} = \textbf{premium} \times ((p_{\max}/\pstar)^2 - 1)$ & \textbf{18.8\%} & $\mathbf{+23.14}$ \\
Exit $\times\; p^3$ ($\alpha=3$) & $\text{payout} = \text{premium} \times ((p_{\max}/\pstar)^3 - 1)$ & 12.1\% & $-14.52$ \\
\bottomrule
\end{tabular}
\end{center}
(All results at $R_0 = \$200\text{K}$ seed capital, $\premfactor = 3.30\%$, $\DeltaStar = 0.09$.)

\begin{definition}[Power-Squared Exit Payoff]\label{def:exit-payoff}
At position exit (\texttt{afterRemoveLiquidity}), PLP $i$ receives:
\begin{equation}\label{eq:exit-payoff}
  \text{payout}_i = \text{premium}_i \times \left[\left(\frac{\max_t\, p(t)}{\pstar}\right)^2 - 1\right]^{+}
\end{equation}
where $\text{premium}_i = \premfactor \times \text{lifetime\_fees}_i$,
$p(t) = \DeltaPlus(t)/(1-\DeltaPlus(t))$ is the concentration price at time $t$,
and $\max_t$ is taken over the position's lifetime.
\end{definition}

\subsection{Why Power-Squared Dominates}

\begin{enumerate}
  \item \textbf{No moral hazard.} The payout depends on the \emph{worst} (highest)
    concentration experienced during the position's lifetime. Staying in the pool
    during concentration spikes \emph{increases} the payout. Exiting early reduces it.
    This is the opposite of the premium-proportional moral hazard.

  \item \textbf{Tail sensitivity.} The $\alpha = 2$ exponent is the sweet spot for
    rare severe events at $\DeltaStar = 0.09$: sufficient convexity to reward
    extreme concentration without draining the reserve at moderate $\DeltaPlus$.
    Higher $\alpha$ (e.g., $\alpha = 3$) over-concentrates payouts on the single worst event;
    lower $\alpha$ (e.g., $\alpha = 1$) provides insufficient tail compensation.

  \item \textbf{Incentive alignment.} Backtest confirms positive correlation between
    position lifetime and hedge value: longer-lived positions benefit more than
    short-lived ones (the correct incentive structure for insurance).

  \item \textbf{Bootstrapping resolution.} Equal pro-rata (INS-05) yields 0\% of
    positions better off at any $\premfactor$ due to the bootstrapping problem:
    the mutual reserve is too small at trigger time. The exit payoff sidesteps
    bootstrapping entirely by settling at exit, not at trigger.
\end{enumerate}

\section{Architectural Design}\label{sec:architecture}

The architecture is a \textbf{Lendgine-style CFMM} following
Leifke \& Scott \cite{leifke2024powermaker}, adapted from ETH price to
fee concentration price $p = \DeltaPlus/(1-\DeltaPlus)$.

\subsection{Participants and Token Mapping}

\begin{definition}[Insurance CFMM Roles]\label{def:cfmm-roles}
\begin{itemize}
  \item \textbf{Underwriter (LP of insurance CFMM):} External actor who deposits USDC
    collateral into a tick range $[\pl, \pu]$. Earns borrow rate (premium stream).
    Commits to $-p^2$ payoff: loses collateral quadratically as concentration rises.

  \item \textbf{PLP (Borrower):} Uniswap V4 LP who borrows the insurance CFMM's LP
    share. Receives $+p^2$ protection payoff. Pays borrow rate from V4 accrued fees.

  \item \textbf{Token X (virtual):} Concentration exposure token---not a real ERC-20
    but an internal accounting unit representing units of concentration risk priced
    by the FCI oracle.

  \item \textbf{Token Y (real):} USDC---underwriter's collateral, PLP's premium currency.
\end{itemize}
\end{definition}

\subsection{LP Payoff (Underwriter)}

Within tick range $[\pl, \pu]$, the underwriter's portfolio value per unit liquidity:
\begin{equation}\label{eq:lp-payoff}
  V(p) = \begin{cases}
    \dfrac{\pu^2 - \pl^2}{\pl^2} & p \leq \pl \\[6pt]
    \dfrac{\pu^2 - p^2}{\pl^2}   & \pl \leq p \leq \pu \\[6pt]
    0                              & p \geq \pu
  \end{cases}
\end{equation}

\begin{proposition}[Concavity]\label{prop:lp-concavity}
$V(p)$ is concave: $V''(p) = -2/\pl^2 < 0$ in $[\pl, \pu]$.
\end{proposition}

\subsection{Borrower Payoff (PLP)}

The PLP borrows the LP share and receives the convex complement:
\begin{equation}\label{eq:plp-payoff}
  N(p) = V(\pl) - V(p) = \begin{cases}
    0                                                    & p \leq \pl \\[6pt]
    \left(\dfrac{p}{\pl}\right)^2 - 1                   & \pl \leq p \leq \pu \\[6pt]
    \left(\dfrac{\pu}{\pl}\right)^2 - 1                 & p \geq \pu
  \end{cases}
\end{equation}
This matches the backtest-confirmed payoff (\cref{def:exit-payoff}) with
$\pstar = \pl$ (the lowest active underwriter tick) and cap at $\pu$.

\subsection{Market-Implied Strike}

The strike $\pstar$ is \emph{not} a protocol parameter. It emerges from the market:
\begin{equation}\label{eq:market-strike}
  \pstar = \min\{\pl : \exists\; \text{underwriter with liquidity at } \pl\}
\end{equation}
Underwriters express their view on where concentration becomes dangerous by
choosing their tick range $[\pl, \pu]$. The aggregate liquidity profile across
ticks is the insurance pricing curve. As market conditions change, underwriters
reposition, and $\pstar$ adjusts naturally---no governance required.

\subsection{Connection to PowerMaker}

The architecture maps directly onto the PowerMaker \cite{leifke2024powermaker}
Lendgine construction:

\begin{center}
\begin{tabular}{@{}lll@{}}
\toprule
PowerMaker & ThetaSwap Insurance & Role \\
\midrule
ETH price $p$ & Concentration price $p = \DeltaPlus/(1-\DeltaPlus)$ & Underlying \\
LP (deposits into CFMM) & Underwriter (deposits USDC) & Sells convexity \\
Borrower (borrows LP share) & PLP (borrows LP share) & Buys convexity \\
$V(p) = 2pk - p^2$ & $V(p) = (\pu^2 - p^2)/\pl^2$ & LP payoff \\
$\varphi = R_1 - (k - \tfrac{1}{2}R_2)^2$ & $\psi = y - \tfrac{\pl^2}{4}u$ & Trading function \\
JumpRate borrow curve & JumpRate + $\premfactor$ base & Borrow rate \\
\bottomrule
\end{tabular}
\end{center}
